\chapter{Conclusiones y Líneas Futuras}
\label{cap:conclusiones}

\section{El Hilo Conductor: del Colapso a su Resolución}
\label{sec:cumplimiento_objetivos}

El Capítulo~\ref{cap:intro} planteaba como hilo conductor de esta memoria la resolución del problema más persistente de este paradigma: el \textbf{representation collapse} (colapso de representación), donde la red minimiza la energía a cero mapeando todas las entradas a una constante trivial $f_{\theta}(x) = c$. Cada capítulo abordó una etapa distinta de ese problema, y conviene cerrar la memoria recorriendo el mismo orden para mostrar cómo cada etapa queda finalmente integrada en el marco de este capítulo.

El Capítulo 2 introdujo los EBMs como marco unificador para el aprendizaje de compatibilidades y sentó, a través del Principio de Máxima Entropía de Jaynes, la base física de toda asignación de probabilidades sobre $\mathcal{Z}$. Ese mismo principio reaparece aquí en su forma más concreta: la distribución gaussiana isótropa que SIGReg fuerza sobre $\mathcal{Z}$ es precisamente la distribución de máxima entropía bajo varianza fija, de modo que el marco de Jaynes introducido en el Capítulo 2 no solo fundamenta la teoría, sino que determina la elección de diseño de este capítulo final.

El Capítulo 3 demostró analíticamente la vulnerabilidad al colapso de representación de las redes siamesas cuando solo se dispone de alineamiento positivo. El Capítulo 4 mostró que el aprendizaje contrastivo evita ese colapso mediante fuerzas de repulsión basadas en muestras negativas, a costa de los límites computacionales estudiados en esa misma sección. SIGReg no pertenece a ninguna de estas dos familias: el colapso queda excluido por construcción, sin pares negativos, porque la regularización exige una varianza fija y una covarianza nula en cada dimensión de $\mathcal{Z}$.

El Capítulo 5 estudió la revolución no contrastiva inspirada en la hipótesis de Barlow, mostrando que Barlow Twins y VICReg evitan el colapso mediante regularizaciones matriciales, pero a cambio de una dependencia que ninguna de las dos arquitecturas reconoce de forma explícita: su garantía de independencia estadística solo es válida bajo normalidad de $Z$. El Capítulo 6 extendió el marco a los world models temporales y a las JEPA, mostrando que PLDM hereda esa misma dependencia gaussiana no verificada, mientras que V-JEPA evita el colapso mediante heurísticas arquitectónicas, EMA y stop-gradient, que funcionan en la práctica sin ningún principio estadístico que las respalde.

El Capítulo 7 expuso el formalismo de LeWorldModel y su simplificación probabilística de regularización isótropa, SIGReg, que utiliza el Teorema de Cramér-Wold y el test de Epps-Pulley para lograr una regularización libre de colapso. Con ello cierra el hilo conductor resolviendo la limitación teórica subyacente: fuerza, en lugar de asumir, la normalidad multivariante que Barlow Twins, VICReg y PLDM necesitaban sin garantizar para que la decorrelación implique independencia. El colapso de representación que abre la memoria en el Capítulo 3 queda así descartado en el Capítulo 7, no por una solución arquitectónica como en V-JEPA, sino por una condición estadística, rigurosa y verificable.

\section{Líneas Futuras}
\label{sec:lineas_futuras}

Esta memoria establece las bases teóricas de los \emph{world models} basados en \emph{Joint Embedding}. Su aplicación práctica, sin embargo, apenas comienza a explorarse, y las limitaciones discutidas a lo largo de este trabajo trazan dos grandes líneas de investigación y desarrollo futuras, alineadas con la visión contemporánea de H-JEPA (\emph{Hierarchical JEPA}) y su extensión a sistemas complejos.

La primera línea, y la más apremiante en el ámbito de la robótica y el control de agentes autónomos, es la transición hacia \textbf{Modelos de Mundo Jerárquicos (H-JEPA)}. Actualmente, los modelos de tipo VLA (\emph{Vision-Language-Action}) son la última tecnología en control robótico; sin embargo, dependen de la imitación del comportamiento (\emph{behavioral cloning}) y carecen de capacidad de planificación explícita al no poseer un modelo predictivo de sus propias acciones. El enfoque JEPA supera esta carencia al aprender un modelo del mundo condicionado por la acción que permite evaluar secuencias y planificar en el espacio latente. No obstante, los modelos predictivos con arquitectura plana, como LeWorldModel, se muestran inestables en horizontes de planificación largos (fallando al intentar predecir más allá de unos pocos pasos en tareas físicas simples como el entorno \emph{Push-T}~\cite{klindt2026lejepa}). Una vía prometedora hacia la solución de este problema es la predicción jerárquica, es decir, estructurar el espacio latente en múltiples niveles de abstracción. En H-JEPA, los niveles inferiores realizarían predicciones a muy corto plazo reteniendo un alto grado de detalle (necesario, por ejemplo, para el control motor de milisegundo a milisegundo), mientras que los niveles superiores operarían sobre horizontes temporales más amplios omitiendo los detalles finos e impredecibles. Las predicciones en la capa superior actúan entonces como sub-objetivos (\emph{subgoals}) para el planificador del nivel inferior. Integrar garantías probabilísticas estables como SIGReg a través de estas múltiples capas para evitar el colapso dimensional es un problema abierto y fundamental para alcanzar un verdadero razonamiento a largo plazo.

La segunda línea, no menos ambiciosa, es la aplicación de estos modelos a \textbf{sistemas dinámicos complejos}. El control óptimo clásico requiere que la dinámica del sistema pueda describirse con unas pocas ecuaciones matemáticas (como la cinemática de un brazo robótico o la dinámica de una población). Sin embargo, arquitecturas como JEPA abren la puerta a aprender simuladores latentes puramente a partir de datos observacionales en dominios donde el modelado analítico resulta intratable. Esto expande el uso de los modelos de mundo más allá de los agentes físicos y virtuales tradicionales, dirigiéndose hacia sistemas industriales como la predicción aerodinámica en motores a reacción, el control de plantas químicas o de energía, el diseño en ciencia de materiales o incluso modelos económicos y sociales complejos. A largo plazo, abarcaría incluso dominios biológicos y médicos, como planificar un curso de tratamiento personalizado simulando el estado fenotípico futuro de un paciente crónico frente a distintas intervenciones terapéuticas. En todos estos escenarios de alta complejidad predictiva, el uso de regularizadores probabilísticos para imponer una estructura estadística bien condicionada en el espacio latente (\emph{embedding space}) resulta ser un paso indispensable.

Ambas líneas convergen, en última instancia, en un mismo problema: dotar al espacio de representación de una estructura estadística que sea a la vez rica y estable, capaz de sostener la predicción sin colapsar bajo su propio peso. Si el cálculo diferencial proporcionó el lenguaje para describir los sistemas continuos de la física clásica, quizás la geometría estadística de los espacios latentes esté llamada a desempeñar un papel semejante para las máquinas que, en las próximas décadas, aprenderán a anticipar el futuro antes de actuar sobre él. Este trabajo no pretende agotar esa tarea, sino señalar, con el rigor que permite el estado actual de la teoría, una de sus primeras rutas transitables.
